\chapter{Conjuntos}
\label{cap:conjuntos}

\section{Conjuntos, inclusión y predicados}
\label{sec:conjuntos-predicados}

Un conjunto está determinado por sus \textbf{elementos}: debe quedar
claro si cada objeto pertenece a él. Se escribe $x\in A$ para indicar
pertenencia y $x\notin A$ para negarla. Conjunto y pertenencia se toman
como nociones primitivas. Un conjunto puede ser elemento de otro;
en la teoría considerada, no es elemento de sí mismo.

Se puede describir un conjunto \textbf{por extensión}, enumerando sus
elementos entre llaves, o \textbf{por comprensión}, mediante una propiedad:
\[
 A=\{x\in U\mid P_x\}.
\]
Aquí $U$ es el \textbf{universo} y $P_x$ un predicado, que se convierte
en proposición al fijar $x\in U$. El conjunto $A$ es la
\textbf{extensión del predicado}, y $P$ una \textbf{propiedad característica}
de $A$ en $U$. Dos predicados son equivalentes en un mismo universo si
seleccionan los mismos elementos. Las conectivas del
capítulo~\ref{cap:logica} permiten formar nuevos predicados.

El orden y las repeticiones en una enumeración no cambian el conjunto.
El \textbf{conjunto unitario} $\{a\}$ tiene un único elemento; hay que
distinguir $a$ de $\{a\}$. El \textbf{conjunto vacío} $\varnothing$ no
contiene elementos y está incluido en cualquier conjunto.
En cambio, $\{\varnothing\}$ tiene un elemento.

La \textbf{inclusión} $A\subseteq B$ significa que todo elemento de $A$
pertenece a $B$; admite $A=B$. La inclusión estricta $A\subsetneq B$
añade $A\ne B$. La igualdad se comprueba por doble inclusión:
\[
 A=B\eqv\bigl(A\subseteq B\ \text{y}\ B\subseteq A\bigr).
\]
Para demostrar $A\subseteq B$, toma un elemento arbitrario de $A$ y
justifica que pertenece a $B$. Los diagramas de Venn representan los
conjuntos por regiones y ayudan a visualizar pertenencia e inclusión.

El \textbf{conjunto de las partes} o conjunto potencia reúne todos los
subconjuntos de $A$:
\[
 \mathcal P(A)=\{B\mid B\subseteq A\}.
\]
Siempre $\varnothing\in\mathcal P(A)$ y $A\in\mathcal P(A)$;
en particular, $\mathcal P(\varnothing)=\{\varnothing\}$.
Si $A$ tiene $n$ elementos, $\mathcal P(A)$ tiene $2^n$: para cada
elemento hay dos posibilidades, incluirlo o no en el subconjunto.
Así, $B\in\mathcal P(A)$ equivale a $B\subseteq A$, pero
pertenencia e inclusión son relaciones distintas.

\newpage
\section{Números naturales, inducción y recurrencia}
\label{sec:induccion}

Se usa la convención $\N=\{0,1,2,\ldots\}$ y
$\N^*=\{1,2,3,\ldots\}$. La presentación informal de los
\textbf{axiomas de Peano} establece:
\begin{enumerate}
 \item $0$ es natural.
 \item Cada natural $n$ tiene un único sucesor $n'$, también natural.
 \item $0$ no es sucesor de ningún natural.
 \item Si dos sucesores son iguales, sus predecesores son iguales.
 \item Un subconjunto de $\N$ que contiene $0$ y el sucesor de cada uno
       de sus elementos contiene todos los naturales.
\end{enumerate}
Los sucesores se denotan habitualmente por $n+1$. El proceso de contar
no termina: $\N$ es infinito. Un conjunto es finito, en sentido intuitivo,
si al contar sus distintos elementos se termina.

El último axioma fundamenta la \textbf{demostración por inducción}.
Para probar $P_n$ para todo $n\in\N$:
\begin{enumerate}
 \item \textbf{Caso inicial:} demuestra $P_0$.
 \item \textbf{Paso inductivo:} para un $n$ arbitrario, supone $P_n$
       y demuestra $P_{n+1}$.
\end{enumerate}
El conjunto de naturales que satisfacen $P$ contiene $0$ y está cerrado
por sucesores; por eso coincide con $\N$. La hipótesis inductiva se
utiliza para demostrar el paso, sin suponer de antemano la conclusión.

Si se quiere demostrar $P_n$ solo para $n\ge a$, con $a\in\N$, el caso
inicial es $P_a$ y el paso se prueba para todo $n\ge a$.
En la \textbf{inducción completa} se demuestra $P_0$ y, para cada $n$,
se deduce $P_{n+1}$ suponiendo ciertas todas las proposiciones
$P_0,\ldots,P_n$. Esta variante permite usar varios casos anteriores.

Una \textbf{definición recurrente} determina un término mediante otros
anteriores. Necesita una regla y suficientes valores iniciales para
determinar los términos sucesivos. Por ejemplo, el factorial queda
definido por
\[
 0!=1,\qquad (n+1)!=(n+1)n!\quad(n\in\N).
\]
Para una progresión, fijado $x_1$, las reglas y sus expresiones son
\[
 \begin{array}{lll}
 \text{Aritmética:}&x_n=x_{n-1}+d,&x_n=x_1+(n-1)d,\\[2pt]
 \text{Geométrica:}&x_n=r x_{n-1},&x_n=x_1r^{\,n-1}.
 \end{array}
\]
Las reglas se aplican para $n\ge2$; el valor de $x_1$ se da por separado.

\newpage
\section{Cuantificadores}
\label{sec:cuantificadores}

Para un predicado $P_x$ sobre $U$, los cuantificadores forman proposiciones:
\begin{itemize}
 \item $(\forall x\in U)\,P_x$: \textbf{universal}, todos los elementos
       de $U$ satisfacen $P$. Equivale a $\{x\in U\mid P_x\}=U$.
 \item $(\exists x\in U)\,P_x$: \textbf{existencial}, al menos un
       elemento de $U$ satisface $P$. Equivale a
       $\{x\in U\mid P_x\}\ne\varnothing$.
\end{itemize}
El nombre de la variable cuantificada puede cambiarse conservando su
universo, cambiando todas sus apariciones y sin confundirla con otra
variable. El universo se omite únicamente cuando está claro.
Para $U=\{u_1,\ldots,u_m\}$, no vacío, el universal equivale a
$P_{u_1}\wedge\cdots\wedge P_{u_m}$ y el existencial a
$P_{u_1}\vee\cdots\vee P_{u_m}$.

\subsection*{Negación y comprobación}

Negar cambia el cuantificador y niega el predicado:
\begin{align*}
 \neg\bigl[(\forall x\in U)P_x\bigr]
   &\eqv(\exists x\in U)\neg P_x,\\
 \neg\bigl[(\exists x\in U)P_x\bigr]
   &\eqv(\forall x\in U)\neg P_x.
\end{align*}
Una proposición universal se refuta con un \textbf{contraejemplo}:
un elemento del universo que no cumpla el predicado. Una existencial
se demuestra dando un elemento que lo cumpla; para refutarla hay que
justificar que ninguno lo cumple.

\subsection*{Varias variables y orden de los cuantificadores}

Un predicado $R_{xy}$ sobre $A\times B$ se convierte en proposición
cuando se cuantifican ambas variables. En
$(\forall x\in A)(\exists y\in B)R_{xy}$, el elemento $y$ puede
depender de $x$. En $(\exists y\in B)(\forall x\in A)R_{xy}$, un mismo
$y$ debe servir para todos los $x$. Estas afirmaciones no son
equivalentes en general.

\begin{ejemplo}
En $\N$, $(\forall x)(\exists y)\,y=x+1$ es verdadero: se elige el
sucesor de cada $x$. En cambio, $(\exists y)(\forall x)\,y=x+1$ es
falso: los casos $x=0$ y $x=1$ exigirían a la vez $y=1$ e $y=2$.
\end{ejemplo}

Para negar varios cuantificadores, se recorren desde fuera hacia dentro,
cambiando cada $\forall$ por $\exists$ y cada $\exists$ por $\forall$,
sin alterar su orden, y se niega el predicado final. Por ejemplo,
\[
 \neg\bigl[(\exists x\in A)(\forall y\in B)R_{xy}\bigr]
 \eqv(\forall x\in A)(\exists y\in B)\neg R_{xy}.
\]

\section{Operaciones y álgebra de conjuntos}
\label{sec:algebra-conjuntos}

Para $A,B\subseteq U$, el \textbf{complementario} se toma siempre
respecto del mismo universo $U$. Las operaciones básicas se definen por:
\begin{align*}
 A\cup B&=\{x\in U\mid x\in A\vee x\in B\},\\
 A\cap B&=\{x\in U\mid x\in A\wedge x\in B\},\\
 \overline A&=\{x\in U\mid x\notin A\}=U\setminus A,\\
 A\setminus B&=\{x\in U\mid x\in A\wedge x\notin B\}
                =A\cap\overline B.
\end{align*}
La \textbf{unión} reúne los elementos que están en al menos uno de los
conjuntos; la \textbf{intersección}, los comunes. Los conjuntos son
\textbf{disjuntos} cuando $A\cap B=\varnothing$.
La \textbf{diferencia} conserva los elementos de $A$ que no están en $B$;
en general, $A\setminus B\ne B\setminus A$.
Cambiar el universo puede cambiar el complementario.

La \textbf{diferencia simétrica} reúne los elementos que pertenecen
exactamente a uno de los dos conjuntos:
\[
 A\mathbin{\triangle}B
  =(A\setminus B)\cup(B\setminus A)
  =(A\cup B)\setminus(A\cap B).
\]
Corresponde a la disyunción exclusiva de los predicados de pertenencia.

\subsection*{Leyes fundamentales}

Para $A,B,C\subseteq U$, se tienen las siguientes identidades:
\begin{center}
\small
\begin{tabular}{@{}p{0.25\linewidth}p{0.69\linewidth}@{}}
\toprule
Ley & Identidades\\
\midrule
Idempotencia & $A\cup A=A$, \quad $A\cap A=A$.\\[3pt]
Conmutatividad & $A\cup B=B\cup A$, \quad $A\cap B=B\cap A$.\\[3pt]
Asociatividad & $(A\cup B)\cup C=A\cup(B\cup C)$,\newline
               $(A\cap B)\cap C=A\cap(B\cap C)$.\\[3pt]
Distributividad & $A\cup(B\cap C)=(A\cup B)\cap(A\cup C)$,\newline
                 $A\cap(B\cup C)=(A\cap B)\cup(A\cap C)$.\\[3pt]
Vacío y universo & $A\cup\varnothing=A$, \quad $A\cap\varnothing=\varnothing$,\newline
                  $A\cup U=U$, \quad $A\cap U=A$.\\[3pt]
Complementario & $A\cup\overline A=U$, \quad $A\cap\overline A=\varnothing$,\newline
                $\overline{\overline A}=A$, \quad
                $\overline U=\varnothing$, \quad $\overline\varnothing=U$.\\[3pt]
De Morgan & $\overline{A\cup B}=\overline A\cap\overline B$,\newline
            $\overline{A\cap B}=\overline A\cup\overline B$.\\
\bottomrule
\end{tabular}
\end{center}
Estas leyes reflejan las equivalencias lógicas de
\S\ref{sec:equivalencias}: pertenecer a una unión corresponde a $\vee$;
a una intersección, a $\wedge$; y a un complementario, a $\neg$.
Se pueden justificar comparando los predicados de pertenencia o
demostrando las dos inclusiones.

La inclusión también se reconoce mediante operaciones:
\[
 A\subseteq B\eqv A\cup B=B\eqv A\cap B=A.
\]
Siempre se cumplen las dos cadenas de inclusiones
\[
 A\cap B\subseteq A\subseteq A\cup B,\qquad
 A\cap B\subseteq B\subseteq A\cup B.
\]

\subsection*{Familias de conjuntos}

Una \textbf{familia de conjuntos} $(F_i)_{i\in I}$ asocia a cada índice
$i$ un conjunto $F_i$; aquí se supone $I\ne\varnothing$.
Si todos los $F_i$ están incluidos en $U$, se definen
\begin{align*}
 \bigcup_{i\in I}F_i
   &=\{x\in U\mid(\exists i\in I)\ x\in F_i\},\\
 \bigcap_{i\in I}F_i
   &=\{x\in U\mid(\forall i\in I)\ x\in F_i\}.
\end{align*}
Para pertenecer a la unión basta estar en un miembro de la familia;
para pertenecer a la intersección hay que estar en todos.
La asociatividad permite escribir las uniones e intersecciones finitas
sin paréntesis.

\newpage
\section{Producto cartesiano}
\label{sec:producto-cartesiano}

Un \textbf{par ordenado} $(x,y)$ distingue primera y segunda componente.
Su igualdad se comprueba componente a componente:
\[
 (x,y)=(u,v)\eqv(x=u\ \text{y}\ y=v).
\]
No es lo mismo que $\{x,y\}$: intercambiar las componentes puede cambiar
el par, y se admite $x=y$.
El \textbf{producto cartesiano} es
\[
 A\times B=\{(x,y)\mid x\in A,\ y\in B\}.
\]
En general, $A\times B\ne B\times A$. Además,
\[
 A\times B=\varnothing\eqv
   (A=\varnothing\ \text{o}\ B=\varnothing).
\]
Si $A'\subseteq A$ y $B'\subseteq B$, entonces
$A'\times B'\subseteq A\times B$.

El producto distribuye sobre unión e intersección en cada factor:
\begin{align*}
 A\times(B\cup C)&=(A\times B)\cup(A\times C),\\
 A\times(B\cap C)&=(A\times B)\cap(A\times C),\\
 (A\cup B)\times C&=(A\times C)\cup(B\times C),\\
 (A\cap B)\times C&=(A\times C)\cap(B\times C).
\end{align*}
Para comprobarlas se descompone la pertenencia de un par en las
condiciones sobre cada componente.

El producto de $n$ conjuntos, para $n\ge1$, reúne las
\textbf{$n$-uplas ordenadas}:
\[
 A_1\times\cdots\times A_n
 =\{(x_1,\ldots,x_n)\mid x_i\in A_i\text{ para }i=1,\ldots,n\}.
\]
Las $n$-uplas son iguales si coinciden todas sus componentes.
Si los factores son el mismo conjunto $A$, se escribe $A^n$.
El plano y el espacio euclídeos se representan mediante $\R^2$ y $\R^3$.

\section{Relaciones entre conjuntos}
\label{sec:relaciones-conjuntos}

Una \textbf{relación} o \textbf{correspondencia} de $A$ a $B$ es un
subconjunto $\mathcal R\subseteq A\times B$. $A$ es el conjunto inicial
y $B$, el final; si $A=B$, se habla de una relación en $A$.
La notación $x\mathcal R y$ significa $(x,y)\in\mathcal R$.
El conjunto de todas las relaciones de $A$ a $B$ es $\mathcal P(A\times B)$.

Un predicado $R_{xy}$ sobre $A\times B$ determina su \textbf{grafo},
\[
 \mathcal R=\{(x,y)\in A\times B\mid R_{xy}\}.
\]
Recíprocamente, pertenecer a un subconjunto de $A\times B$ define un
predicado de dos argumentos. Una relación se puede describir por
extensión o comprensión, y representar mediante puntos o flechas.

La \textbf{relación inversa} intercambia las componentes:
\[
 \mathcal R^{-1}
   =\{(y,x)\in B\times A\mid(x,y)\in\mathcal R\}.
\]
\begin{samepage}
Para cada elemento se consideran los siguientes conjuntos:
\begin{align*}
 \mathcal R(x)&=\{y\in B\mid(x,y)\in\mathcal R\}
          &&\text{(imagen de $x\in A$)},\\*
 \mathcal R^{-1}(y)&=\{x\in A\mid(x,y)\in\mathcal R\}
          &&\text{(original de $y\in B$)}.
\end{align*}
\end{samepage}
\begin{samepage}
La \textbf{imagen de la relación} y su \textbf{conjunto original} son
\begin{align*}
 \mathcal R(A)&=\{y\in B\mid(\exists x\in A)\ (x,y)\in\mathcal R\},\\*
 \mathcal R^{-1}(B)&=\{x\in A\mid(\exists y\in B)\ (x,y)\in\mathcal R\}.
\end{align*}
\end{samepage}
No tienen por qué coincidir con $B$ y $A$, respectivamente: puede haber
elementos sin pareja. La imagen o el original de un elemento es un
\emph{conjunto}, que puede ser vacío o contener varios elementos.

\begin{ejemplo}
Para $A=\{a,b,c\}$, $B=\{1,2,3\}$ y
$\mathcal R=\{(a,1),(a,2),(b,2)\}$,
\[
 \mathcal R(a)=\{1,2\},\quad \mathcal R(c)=\varnothing,\quad
 \mathcal R^{-1}(2)=\{a,b\}.
\]
Además, $\mathcal R(A)=\{1,2\}$ y $\mathcal R^{-1}(B)=\{a,b\}$.
\end{ejemplo}

\subsection*{Composición}

Si $\mathcal R\subseteq A\times B$ y $\mathcal S\subseteq B\times C$,
su \textbf{composición} es la relación de $A$ a $C$ definida por
\[
 \mathcal S\circ\mathcal R
 =\{(x,z)\in A\times C\mid
       (\exists y\in B)\bigl(x\mathcal R y\wedge y\mathcal S z\bigr)\}.
\]
Para calcularla, se buscan pasos de $x$ a un elemento intermedio $y$
mediante $\mathcal R$ y de ese mismo $y$ a $z$ mediante $\mathcal S$.
Basta un intermediario; si varios producen el mismo par, se incluye
una sola vez. En $\mathcal S\circ\mathcal R$ actúa primero
$\mathcal R$.

Los predicados de varias variables se combinan con las mismas
conectivas y cuantificadores. Su interpretación y negación siguen las
reglas de \S\ref{sec:cuantificadores}, respetando el universo de cada variable.

\section{Límites de la descripción intuitiva}

Una propiedad no permite formar sin restricciones un conjunto de
\emph{todos} los objetos que la cumplen. La comprensión usada en
\S\ref{sec:conjuntos-predicados} selecciona elementos dentro de un
conjunto $U$ ya dado. La colección de todos los conjuntos no constituye
un conjunto.

La \textbf{paradoja de Russell} muestra la dificultad: si la colección
$M$ de todos los conjuntos que no se pertenecen a sí mismos fuera un
conjunto, se tendría
\[
 M\in M\eqv M\notin M,
\]
lo que es contradictorio.

Una teoría axiomática, como Zermelo--Fraenkel (ZF), precisa qué
construcciones de conjuntos se admiten. La igualdad se rige por sus
elementos (extensionalidad), y la especificación permite seleccionar
mediante una propiedad los elementos de un conjunto existente.
Los axiomas sustituyen la formación irrestricta de colecciones y
sirven de base para justificar los resultados.
